% !TEX root = ../print.tex
% Section 5.7 of the article: the overview table and the inclusions (notes/PLAN-article-structure.md,
% "The outline"). The table and the inclusions are prose of the article; the one node here is
% thm:wedge-axis, the general linear row along the chain of invariant cones. It moved from the end of
% 05-wedges.tex (5.3) on the author's decision of 2026-10-04, so that it follows every node it \uses
% (the stable rung and the ordered-density cone of 5.5, the budget and the stable-order cell of 5.6).
% No label, statement or proof changed in the move (scripts/check-statements-unchanged.py).

\section{The overview table and the inclusions}
\label{sec:landscape-table}

The overview table has two axes, the covariance group and the wedge. Along the general linear row
the families are read off the chain of invariant cones of
\ref{sec:invariant-wedges}: the larger the cone, the fewer the families, and the stable cones fill the
chain below the Gaussian's. The row's other cells, and the rows of the smaller groups, are the nodes of
\ref{sec:affine-orbits} to \ref{sec:smaller-groups}.

% not in the draft: the plan's item D3, the chain rays < K_alpha < K_beta < PSD with the gap at 2 (Q-0266).
% What writing it down found: along a cone-generated order inside PSD(d) and away from its boundary, one
% number of the cone, beta(C), the least alpha with C inside C_alpha, decides which stable orbits are
% families and whether non-Gaussian families exist; the gap at 2 has two mechanisms, the
% rank-deficient step of thm:rank-deficient-step when the cone reaches the boundary of PSD(d), and the
% anisotropy budget of prop:anisotropy-budget between C_2 and that boundary; and cones not inside PSD(d)
% carry no family at all.
% CHANGED Q-0281: typed (the statement skeleton). Clauses (1), (2)(a) and (2)(b), which spend
% nothing, are split to lem:wedge-axis-chain, placed after this node (ADR-0025;
% records/formalization/PRICING.md: "(1), (2)(a)-(b) [T] vs (2)(c), (3)"), and \uses gains it; the
% text is unchanged and the \lean tag names the child's declarations too.
% CHANGED Q-0281: \uses re-pointed to the split child it reads (thm:cone-orders(4) beside (1); prop:similarity-stable-order-families(1) beside (2)).
% CHANGED Q-0281 wave 2: entry-free route (the author's decision on W2F-1). \uses re-pointed to what
% the proof reads: thm:cone-orders, lem:cone-orders-psd and prop:invariant-orders dropped (the orders are
% read through prop:stable-order-generated and the causal curves, without ledger A9); the statement
% still names thm:cone-orders(4) for the definition of the order on PSD(d), which spends nothing.
% CHANGED Q-0281 wave 2: blind review W2H-1: \uses gains lem:spectral-cones (the last sentence of its
% (4), which spends nothing, in the proof of the antisymmetry); the assignment says so.
% CHANGED Q-0281 wave 2: blind review W2H-2: \uses gains def:cone-orders-curves, where the vocabulary of
% the statement (C-causal curve, the orders on PD(d) and PSD(d)) now has its own node; the statement
% text, shared with the paper, is unchanged and still names thm:cone-orders(4), which restates it.
% CHANGED Q-0381 (2026-10-10; the fidelity review's paper-facing point 6): the last sentence of
% (2)(c), "if beta(C) >= 2, the families are those of (b)", is written out. (b) is stated for another
% cone, so "those of (b)" could be read as the families along that cone's order; what is meant, and
% what wedge_axis_non_gaussian states (its second conjunct, an iff with EqOnPairs (causalPlusLE C)),
% is the affine Gaussian scale spaces restricted to the pairs of this C's order. No claim changes.
% CHANGED writing-rules pass (2026-10-14; records/affine/PASS-writing-rules-affine.md, owed item 1):
% title only, "the wedge axis of the general linear row" -> "the general linear row along the chain of
% cones"; "wedge axis" is a term of the plan the paper never defines. Label, statement and \lean unchanged.
% CHANGED Q-0469 (wording only; records/affine/GATE-affine.md, Q-0317 row "second half 26"): clause (3),
% about ninety words and three semicolons in one sentence, is a list (a)-(d), one item per semicolon-
% separated claim, the text of each unchanged. No claim changes; the clause keeps its number.
% CHANGED round-3 review (report E, section 3.1; the author's decision D10): the annotation's "Not
% claimed" on the converse of the profile clause and on cones with the same beta(C) is replaced by a
% pointer to prop:wedge-axis-profiles, after lem:wedge-axis-chain, which proves both in print. The
% statement is unchanged, so this node's Lean still covers all of it. Annotation only.
% CHANGED D14 (2026-11-04): the pointer says "for every profile" where it said "for a right-continuous
% profile", prop:wedge-axis-profiles(1) having dropped the hypothesis. Annotation only.
% CHANGED round 3, the Lean of prop:wedge-axis-profiles: the annotation's pointer no longer says that
% that node's Lean is owed. Annotation only.
\begin{theorem}[the general linear row along the chain of cones]
\label{thm:wedge-axis}
\uses{def:cone-orders-curves,prop:stable-order-generated,lem:spectral-cones,def:affine-cascade-family,def:gw-cascade-family,lem:gw-kernels-from-origin,prop:gw-ray-families,thm:rank-deficient-step,thm:affine-rigidity,lem:elliptical-stable,prop:stable-step-cone,prop:orbit-step-cone,prop:similarity-stable-order-families,prop:anisotropy-budget,prop:stable-exponent-cone,lem:stable-order-psd,lem:wedge-axis-chain,thm:cited-A3,thm:cited-A5}
\lean{AffineHemigroup.wedge_axis_chain, AffineHemigroup.wedge_axis_not_psd,
AffineHemigroup.wedge_axis_singular, AffineHemigroup.wedge_axis_index,
AffineHemigroup.wedge_axis_stable, AffineHemigroup.wedge_axis_non_gaussian,
AffineHemigroup.wedge_axis_row, AffineHemigroup.wedge_axis_row_members}\leanok
Let $d \ge 2$, and let $C_\alpha$ and $\preceq_{C_\alpha}$ be as in \ref{prop:stable-order-generated} and
$\preceq^+_C$ as in \ref{thm:cone-orders}(4).
\begin{enumerate}
\item \emph{(The chain.)} For $0 \le \alpha < \beta \le \infty$, $C_\alpha \subsetneq C_\beta$ and
$\preceq_{C_\alpha} \subsetneq \preceq_{C_\beta}$. For $\alpha < \infty$ every nonzero member of $C_\alpha$ is positive
definite, while $C_\infty = \PSD(d)$ contains every nonzero singular positive semidefinite matrix. The union of
the $C_\alpha$, $0 \le \alpha < 2$, is $\{0\} \cup \{X \in \Sym(d) : (d+2)\lambda_{\min}(X) > \tr X\}$, a proper
subset of $C_2$.
\item \emph{(Which cones carry general linear families.)} Let $C \subseteq \Sym(d)$ be an $\Orth(d)$-invariant
closed convex cone with $I \in C$ and $C \cap (-C) = \{0\}$, and consider the
$(\GL(d), \preceq^+_C)$-cascade families of kernels (\ref{def:gw-cascade-family}(2)).
\begin{enumerate}
\item[(a)] If $C \not\subseteq \PSD(d)$, there is none.
\item[(b)] If $C \subseteq \PSD(d)$ and $C$ contains a nonzero singular matrix, they are exactly the
restrictions to the pairs $\Sigma \preceq^+_C \Sigma'$ of the affine Gaussian scale spaces
\ref{eq:affine-gaussian}, $b > 0$.
\item[(c)] Otherwise every nonzero member of $C$ is positive definite, the largest value $\beta(C)$ of
$\tr X/\lambda_{\min}(X) - d$ over $X \in C \setminus \{0\}$ is attained, and $\beta(C) \in [0, \infty)$ is the
least $\alpha$ with $C \subseteq C_\alpha$. For $0 < \gamma < 2$, $\Sigma \mapsto (\Sigma^{1/2})_*S_\gamma(I)$
is the family of kernels from scale zero of a $(\GL(d), \preceq^+_C)$-cascade family if and only if
$\gamma \ge \beta(C)$. A $(\GL(d), \preceq^+_C)$-cascade family whose $\mu_I$ is not Gaussian exists if and only
if $\beta(C) < 2$; if $\beta(C) \ge 2$, the families are, as in (b), exactly the restrictions to the pairs
$\Sigma \preceq^+_C \Sigma'$ of the affine Gaussian scale spaces \ref{eq:affine-gaussian}, $b > 0$.
\end{enumerate}
\item \emph{(The row along the chain.)} For $0 \le \alpha \le \infty$ let $\mathfrak F_\alpha$ be the set of
laws $\mu_I$ of $(\GL(d), \preceq^+_{C_\alpha})$-cascade families of kernels.
\begin{enumerate}
\item[(a)] $\mathfrak F_0$ is the set of rotation-invariant self-decomposable laws other than $\delta_0$.
\item[(b)] $\mathfrak F_\beta \subseteq \mathfrak F_\alpha$ for $\alpha \le \beta$, strictly when
$\alpha < \beta \le 2$.
\item[(c)] For $0 < \alpha < 2$, $\mathfrak F_\alpha$ contains $S_\gamma(I)$, $0 < \gamma < 2$, if and only
if $\gamma \ge \alpha$, and contains the law $\mu_I$ of \ref{prop:orbit-step-cone}(1) for every profile whose
dilation ratio has $R(\lambda) \ge \lambda^\alpha$ for $\lambda \ge 1$.
\item[(d)] For $2 \le \alpha \le \infty$, $\mathfrak F_\alpha$ is the set of centred Gaussian laws with
covariance $2bI$, $b > 0$.
\end{enumerate}
\end{enumerate}
\statusA\quad\ledger{A5}\quad\ledger{A3}\quad\emph{(\textbf{Assignment.} \ledger{A5} and \ledger{A3}(i)--(ii)
are spent through \ref{prop:anisotropy-budget}(1), once, in the necessity half of (2)(c) and so in (3) for
$2 \le \alpha < \infty$, including the strictness at $\beta = 2$; \ledger{A3}(iii) through \ref{lem:elliptical-stable}(2), \ref{prop:stable-step-cone}(1)
and \ref{prop:similarity-stable-order-families}(2), for the stable orbits and the profiles. Clauses (1), (2)(a)
and (2)(b) spend nothing: (b) is \ref{thm:rank-deficient-step}, which cites nothing. For cascade families
along $\preceq^+_C$, which has the path condition, entry A6 would give the infinite divisibility for which the
budget spends A5 (\ref{thm:cascade-infinitely-divisible}(2)); the node does not take that route.
Nor does it spend entry A9, through which \ref{thm:cone-orders}(4) makes $\preceq^+_C$ a partial order:
the proof shows the antisymmetry directly, so that $(\GL(d), \preceq^+_C)$ is a wedge datum and (2) and (3)
have their meaning without it, reading only the last sentence of \ref{lem:spectral-cones}(4), which spends
nothing. The Lean does not need the antisymmetry, as its cascade families are stated for any relation.
\textbf{What it is worth.} It is the organizing statement of the plan's wedge axis. The cones of the module form
a chain $C_0 \subsetneq C_\alpha \subsetneq C_\beta \subsetneq C_2 \subsetneq C_\infty = \PSD(d)$, with
$C_0$ the rays and $C_\alpha = \mathfrak k_\alpha \cap \Sym(d)$ for $0 < \alpha < 2$ (\ref{prop:stable-order-generated}(2);
the full exponent cones are $\mathfrak k_\alpha = C_\alpha \oplus$ the skew matrices). Along it the general linear row
shrinks from every self-decomposable orbit to the Gaussian, the stable index deciding the wedge exactly
for the stable orbits. The gap at $2$ has two causes. When the cone reaches the boundary of $\PSD(d)$, whose
nonzero points are the nonzero singular matrices, a single rank-deficient step $I \to e^X$ forces the
Gaussian (\ref{thm:rank-deficient-step}), and no $C_\alpha$ with $\alpha < \infty$ reaches it. Between $C_2$ and
that boundary the budget $\tr B < (d+2)\lambda_{\min}(\mathrm{sym}\,B)$ of \ref{prop:anisotropy-budget}
excludes every family with a non-Gaussian $\mu_I$, at $\alpha = 2$ itself included, since the stable cones fill only the open part of
$C_2$. For a cone in case (c), the stable orbits that are families and the existence of non-Gaussian
families are decided by one number, $\beta(C)$ (set $\beta(C) := \infty$ in case (b)). \textbf{The
converse of the profile clause of (3)(c)}, and the fact that two cones with the same $\beta(C)$ carry the
same laws $\mu_I$, are \ref{prop:wedge-axis-profiles}: for every profile, with any Gaussian part,
$\mu_I \in \mathfrak F_\alpha$ if and only if $R(\lambda) \ge \lambda^\alpha$ for every $\lambda \ge 1$, and with the
polar criterion (\ref{thm:cited-A2}) this describes $\mathfrak F_\alpha$ completely. They are a separate node, so that this node's Lean covers all of it (ADR-0025). \textbf{Not claimed}: the single-step
converse of \ref{prop:orbit-step-cone}, for a step not embedded in a cascade; orders on $\PD(d)$ that are not generated by a cone
(\ref{thm:cone-orders}(3)); the Loewner order of $\PSD(d)$ itself, whose families are those of
\ref{thm:affine-rigidity} and which is larger than $\preceq^+_{C_\infty}$; and any statement for groups
smaller than $\GL(d)$, where \ref{prop:similarity-stable-order-kernels} has the similarity row along
$\preceq^+_{C_\alpha}$, $0 < \alpha < 2$.)}
\end{theorem}

\begin{proof}
\leanok
% CHANGED Q-0281 wave 2: entry-free route (the author's decision on W2F-1). This paragraph is new: it
% proves what the proof reads of the orders, so that no step reads thm:cone-orders(1), (4) or
% lem:gw-kernels-from-origin(1)-(2) with (GL(d), the cone order) as a wedge datum, whose antisymmetry
% spends ledger A9.
% CHANGED Q-0281 wave 2: blind review W2H-1. The antisymmetry of the cone order on PSD(d), "not used"
% before, is proved in this paragraph without ledger A9, in steps (i)-(iii): it is what makes
% (GL(d), the cone order) a wedge datum, which the statements of (2) and (3) presuppose
% (def:gw-cascade-family), here and in lem:wedge-axis-chain(2). New mathematics in a green node's proof;
% the Lean is unaffected, because IsGWCascadeFamily takes any relation, so the node stays green. Step
% (i) reads the last sentence of lem:spectral-cones(4), whose declaration
% AffineHemigroup.orthInvariant_pointed_cone_sign prints Lean core (CIAxiomGuard.lean); it does not
% reach logSpec_product_majorized, which in SpectralCones.lean only expSet_specCone, clause (5), uses.
% CHANGED Q-0281 wave 2: blind review W2H-3: the restriction also reads the invariance of the smaller
% order under the smaller group.
% CHANGED writing-rules pass (2026-10-14; records/affine/PASS-writing-rules-affine.md, owed item 6):
% the proof takes the wording of the paper's copy (Appendix C), without the ledger accounting: "without
% the entry", "which thm:cone-orders(4) obtains through ledger A9 ... without it", "that sentence spends
% no ledger entry" and "defined without ledger A9" go. Which entries the node spends stays in its
% annotation; no step of the argument changes.
\emph{The orders.} The proof reads of $\preceq_C$ and $\preceq^+_C$ only what the causal
curves give directly. By facts (c) and (a) of the proof of \ref{prop:stable-order-generated}: for $X \in C$
the curve $t \mapsto e^{tX}$ is $C$-causal from $I$ to $e^X$, so $I \preceq_C e^X$; as $(\log c)I \in C$ for
$c \ge 1$, $\Sigma \preceq_C c\Sigma$ for $\Sigma \in \PD(d)$; and if $C \subseteq \PSD(d)$, a $C$-causal curve
$\gamma$ has $\gamma' = \gamma^{1/2}(\gamma^{-1/2}\gamma'\gamma^{-1/2})\gamma^{1/2} \ge 0$ almost everywhere, so
it is Loewner-monotone and $\preceq^+_C$ is contained in the Loewner order. A $C$-causal curve is
$C'$-causal for every cone $C'$ containing $C$, so $\preceq^+_C \subseteq \preceq^+_{C'}$. Two $C$-causal curves
$\gamma_1$, $\gamma_2$ with $\gamma_1(1) = \gamma_2(0)$ run in succession at double speed, $\gamma_1(2t)$ for
$t \le 1/2$ and $\gamma_2(2t-1)$ after, form a $C$-causal curve: its entries are absolutely continuous, each being its value at $0$ plus the
integral of the derivatives of $\gamma_1$ and $\gamma_2$ run at double speed, and its relative derivative is twice that of
$\gamma_1$ or $\gamma_2$, which lies in the cone $C$. So $\preceq_C$ is transitive, and with the rays so is
$\preceq^+_C$. For $A \in \GL(d)$, $A\gamma A^{\tp}$ is $C$-causal with $\gamma$, $C$ being $\Orth(d)$-invariant: its
relative derivative is
$Q(\gamma^{-1/2}\gamma'\gamma^{-1/2})Q^{\tp}$ with $Q := (A\gamma A^{\tp})^{-1/2}A\gamma^{1/2}$ orthogonal. So
$\preceq^+_C$ is reflexive, transitive and $\GL(d)$-invariant on $\PSD(d)$ and contains the rays.
Its antisymmetry is that of \ref{thm:cone-orders}(4), and the causal curves give it directly.
(i) Every nonzero $X \in C$ has $\tr X > 0$, by the last sentence of \ref{lem:spectral-cones}(4), as
$I \in C$. (ii) Let $\gamma$ be $C$-causal. By fact (b) of the proof of
\ref{prop:stable-order-generated}, $\det\gamma$, a polynomial in the absolutely continuous entries of
$\gamma$, is absolutely continuous on $[0,1]$, and at almost every $t$ Jacobi's formula gives
$(\det\gamma)' = \det\gamma\,\tr Y$ with $Y := \gamma^{-1/2}\gamma'\gamma^{-1/2} \in C$, so $(\det\gamma)' \ge 0$
by (i); being the integral of its derivative, $\det\gamma$ is nondecreasing. Now let $\Sigma \preceq_C \Sigma'$
and $\Sigma' \preceq_C \Sigma$, along $C$-causal curves $\gamma_1$ from $\Sigma$ to $\Sigma'$ and $\gamma_2$ from
$\Sigma'$ to $\Sigma$. Then $\det\Sigma \le \det\Sigma' \le \det\Sigma$, so $\det\gamma_1$, nondecreasing with
equal values at $0$ and $1$, is constant, and its derivative $\det\gamma_1\,\tr Y_1$ vanishes almost
everywhere. As $\det\gamma_1 > 0$, $\tr Y_1 = 0$ almost everywhere, and so $Y_1 = 0$ almost everywhere by (i).
Hence $\gamma_1' = \gamma_1^{1/2}Y_1\gamma_1^{1/2} = 0$ almost everywhere, $\gamma_1$ is constant by fact (a),
and $\Sigma = \Sigma'$. (iii) Let $\Sigma \preceq^+_C \Sigma'$ and $\Sigma' \preceq^+_C \Sigma$. Only $0$ lies
below $0$, since $0 = c\,\Sigma''$ with $c \ge 1$ forces $\Sigma'' = 0$ and $0$ is not definite; so if one of
the two is $0$, so is the other. Otherwise both are nonzero, and both are singular or both are definite:
a matrix above a nonzero singular $\Sigma$ is a multiple $c\,\Sigma$, the third case of the definition
needing $\Sigma$ definite, and a matrix above a definite one is definite. If both are singular,
$\Sigma' = c\,\Sigma$ and $\Sigma = c'\Sigma'$ with $c, c' \ge 1$, so $cc' = 1$ and $c = 1$. If both are
definite, each of the two relations is $\preceq_C$, as $\Sigma \preceq_C c\Sigma$ for $c \ge 1$, and (ii) gives
$\Sigma = \Sigma'$. So $\preceq^+_C$ is a partial order, and $(\GL(d), \preceq^+_C)$ is a wedge datum
(\ref{def:gw-cascade-family}), $\GL(d)$ being a closed subgroup containing the dilations. A family along one
order is restricted to a smaller order by the argument of \ref{lem:gw-kernels-from-origin}(1), which reads
of the two orders only that one contains the other and that the smaller is invariant under the smaller
group, which every order restricted to here is; and a function that agrees on the pairs
$\Sigma \preceq^+_C \Sigma'$ with a $(\GL(d), \preceq^+_C)$-cascade family is one, as (K1)--(K4) read it only
on pairs, which $\preceq^+_C$ being transitive and $\GL(d)$-invariant keeps among themselves.

% CHANGED Q-0281 wave 2: entry-free route (the author's decision on W2F-1). The strict inclusion of the
% orders is read through prop:stable-order-generated(1), in place of thm:cone-orders(1) and
% prop:invariant-orders(1).
\emph{(1).} Let $0 \le \alpha < \beta \le \infty$. A positive semidefinite $X$ has $\lambda_{\min}(X) \ge 0$, so
$(d+\beta)\lambda_{\min}(X) \ge (d+\alpha)\lambda_{\min}(X) \ge \tr X$ when $\beta < \infty$, and $C_\alpha
\subseteq C_\beta$ (trivially if $\beta = \infty$). Strictness: for $\beta < \infty$,
$X := \diag(1, \dots, 1, 1 + \beta)$ has $\tr X = d + \beta$ and $\lambda_{\min}(X) = 1$, so it lies in $C_\beta$
and not in $C_\alpha$; for $\beta = \infty$, $X := e_1e_1^{\tp} \in \PSD(d)$ has $\lambda_{\min} = 0 < 1 = \tr$,
so it lies in no $C_\alpha$, $\alpha < \infty$. A $C_\alpha$-causal curve is $C_\beta$-causal, so
$\preceq_{C_\alpha} \subseteq \preceq_{C_\beta}$, and $I \preceq_{C_\beta} e^X$ for the $X$ just found. But
$I \not\preceq_{C_\alpha} e^X$: by \ref{prop:stable-order-generated}(1), $I \preceq_{C_0} e^X$ would make $e^X = cI$,
hence $X = (\log c)I \in C_0$, and for $\alpha > 0$, $I \preceq_{C_\alpha} e^X$ would give
$e^{(d+\alpha)\lambda_{\min}(X)} = \lambda_{\min}(e^X)^{d+\alpha} \ge \det e^X = e^{\tr X}$, hence $X \in C_\alpha$.
So $\preceq_{C_\alpha} \subsetneq \preceq_{C_\beta}$. For $\alpha < \infty$ and $0 \ne X \in C_\alpha$, $\tr X > 0$ and
$\lambda_{\min}(X) \ge \tr X/(d+\alpha) > 0$. The union: a nonzero $X$ in some $C_\alpha$, $\alpha < 2$, has
$\lambda_{\min}(X) > 0$ and so $(d+2)\lambda_{\min}(X) > (d+\alpha)\lambda_{\min}(X) \ge \tr X$; conversely
$(d+2)\lambda_{\min}(X) > \tr X \ge d\lambda_{\min}(X)$ gives $\lambda_{\min}(X) > 0$, so $X \ge 0$, and
$X \in C_\alpha$ with $\alpha := \tr X/\lambda_{\min}(X) - d \in [0, 2)$. The matrix $\diag(1, \dots, 1, 3)$ lies in
$C_2$, with equality, and not in the union.

\emph{(2), the kernels from scale zero.} Let a $(\GL(d), \preceq^+_C)$-cascade family be given. The ray order is
contained in $\preceq^+_C$ (its two cases are two of the three), so the family restricts to a
$(\GL(d), \text{rays})$-family with the same kernels from scale zero, and
\ref{prop:gw-ray-families}(3) gives $\mu_\Sigma = (\Sigma^{1/2})_*\mu_I$ with $\mu_I$ rotation invariant,
self-decomposable and not $\delta_0$, and $\widehat\mu_\Sigma(\omega) = F(\langle\omega,\Sigma\omega\rangle)$ with
$F(u) = \widehat\mu_I(\sqrt u\,e_1)$ continuous, $F(0) = 1$. $F$ has no zero (\ref{prop:gw-ray-families}(1)) and is
real ($\mu_I$ is invariant under $-I$), so $F > 0$, and $g := -\log F$ is continuous, $g(0) = 0$, $g \ge 0$ (as
$|F| \le 1$), and $g \not\equiv 0$ (as $\mu_I \ne \delta_0$). For $X \in C$, $I \preceq^+_C e^X$, and (K1) at
$0 \preceq^+_C I \preceq^+_C e^X$ gives $\mu_{e^X} = \mu_{I,e^X} * \mu_I$, that is
\[
  e^{-g(\langle\omega, e^X\omega\rangle)} = e^{-g(|\omega|^2)}\,\widehat\mu_{I,e^X}(\omega)
  \qquad (\omega \in \RR^d), \qquad (\ast)
\]

\emph{(2)(a).} Let $X \in C$ have an eigenvalue $\lambda < 0$ with unit eigenvector $v$, and $m := e^\lambda < 1$.
At $\omega = uv$, ($\ast$) and $|\widehat\mu_{I,e^X}| \le 1$ give $g(mu^2) \ge g(u^2)$, so $g(x) \le g(mx) \le
\dots \le g(m^nx) \to g(0) = 0$ for every $x \ge 0$, against $g \not\equiv 0$.

% CHANGED Q-0281 wave 2: entry-free route (the author's decision on W2F-1). The increments are read
% from (K1) at 0, Sigma, Sigma' in place of lem:gw-kernels-from-origin(2), and the inclusion in the
% Loewner order from the causal curves in place of thm:cone-orders(4).
\emph{(2)(b).} Let $0 \ne X \in C$ be singular and $P := e^X - I$, positive semidefinite with $\rank P =
\rank X$, since its eigenvalues $e^{x_i} - 1$ vanish exactly where those of $X$ do; so $0 < \rank P < d$. ($\ast$)
is \ref{eq:one-step} with $\Sigma = I$ and $\rho = \mu_{I,I+P}$, and \ref{thm:rank-deficient-step} gives
$g(u) = bu$ with $b = g(1)$, $b > 0$ as $g \not\equiv 0$. So $\widehat\mu_\Sigma(\omega) =
e^{-b\langle\omega,\Sigma\omega\rangle}$ for every $\Sigma$. For $\Sigma \preceq^+_C \Sigma'$, (K1) at
$0 \preceq^+_C \Sigma \preceq^+_C \Sigma'$ gives $\widehat\mu_{\Sigma'} = \widehat\mu_{\Sigma,\Sigma'}\,\widehat\mu_\Sigma$,
and $\widehat\mu_\Sigma$ has no zeros, so $\widehat\mu_{\Sigma,\Sigma'} = e^{-b\langle\omega,(\Sigma'-\Sigma)\omega\rangle}$,
the transform of the affine Gaussian kernel ($\Sigma \le \Sigma'$ in the Loewner order, as
$C \subseteq \PSD(d)$): the family is the restriction of the affine Gaussian scale space. Conversely, the affine
Gaussian scale space with constant $b > 0$ is an affine cascade family (\ref{thm:affine-rigidity}), that is a
$(\GL(d), \text{Loewner})$-family; $\preceq^+_C$ is contained in the Loewner order as $C \subseteq \PSD(d)$, so
its restriction is a $(\GL(d), \preceq^+_C)$-family, and so is every function that agrees with it on the pairs.

% CHANGED Q-0281 wave 2: proof rewritten to the Lean route (the slice tr X = 1 in place of the unit
% sphere of a norm).
\emph{(2)(c), the number $\beta(C)$.} Now $C \subseteq \PSD(d)$ has no nonzero singular member, so its nonzero
members are positive definite. On them $f(X) := \tr X/\lambda_{\min}(X) - d$ is continuous,
nonnegative and invariant under positive scaling. The slice $\{X \in C : \tr X = 1\}$ is closed, and bounded,
since a positive semidefinite $X$ with $\tr X = 1$ has entries in $[-1, 1]$; it contains $I/d$, and every
$X \in C \setminus \{0\}$ has the multiple $X/\tr X$ in it. So $f$ attains on the slice a largest value
$\beta(C) \ge 0$, which is its largest value on $C \setminus \{0\}$. For $\alpha \in [0, \infty)$,
$C \subseteq C_\alpha$ means $(d+\alpha)\lambda_{\min}(X) \ge \tr X$, that is $\alpha \ge f(X)$, for every nonzero
$X \in C$, which is $\alpha \ge \beta(C)$.

% CHANGED Q-0281 wave 2: entry-free route (the author's decision on W2F-1). C contained in C_gamma makes
% every C-causal curve C_gamma-causal, in place of exp C contained in exp C_gamma.
\emph{(2)(c), the stable orbits.} If $\gamma \ge \beta(C)$, then $C \subseteq C_\gamma$, so every $C$-causal
curve is $C_\gamma$-causal and $\preceq^+_C \subseteq \preceq^+_{C_\gamma}$; by \ref{prop:stable-order-generated}(1),
$\preceq^+_{C_\gamma}$ is the order $\le_\gamma$ of \ref{prop:similarity-stable-order-families}(1), and that
proposition's clause (2), read with its $\alpha$ equal to $\gamma$ and the stable law of index $\gamma$, gives a
$(\GL(d), \le_\gamma)$-family with kernels $(\Sigma^{1/2})_*S_\gamma(I)$, which restricts. Conversely, if these
are the kernels of a $(\GL(d), \preceq^+_C)$-family and $X \in C$, then $\mu_{e^X} = (e^{X/2})_*S_\gamma(I) =
S_\gamma(e^X)$ (\ref{lem:elliptical-stable}(2)) is $S_\gamma(I) * \mu_{I,e^X}$, and
\ref{prop:stable-step-cone}(1), (a)$\Rightarrow$(c), with $\Sigma = I$ and $\Sigma' = e^X$, gives
$\lambda_{\min}(e^X)^{d+\gamma} \ge \det e^X$, that is $X \in C_\gamma$ (\ref{prop:stable-order-generated}(1)).
So $C \subseteq C_\gamma$ and $\gamma \ge \beta(C)$.

% CHANGED Q-0281 wave 2: proof rewritten to the Lean route (the Gaussian end: g read on the
% one-dimensional marginal, in place of the covariance of a rotation-invariant Gaussian law).
\emph{(2)(c), the non-Gaussian families.} If $\beta(C) < 2$, take any index $\gamma$ with
$\beta(C) \le \gamma < 2$ and $\gamma > 0$, for instance $\gamma := \max(\beta(C), 1)$. The
stable orbit of index $\gamma$ is a family, and $S_\gamma(I)$ is not Gaussian, its transform
$e^{-|\omega|^\gamma}$ having modulus $e^{-t^\gamma}$ at $\omega = te_1$, which is $e^{-at^2}$ for no $a \ge 0$.
Conversely let a family have a non-Gaussian $\mu_I$, and let $0 \ne X \in C$ and $t \ge 0$. By (K1), as in ($\ast$) with
$tX \in C$ in place of $X$, and the form of the kernels, $(e^{tX/2})_*\mu_I = \mu_{e^{tX}} = \mu_I * \mu_{I,e^{tX}}$; applying
$(e^{-tX/2})_*$, $\mu_I = (e^{-tB})_*\mu_I * (e^{-tB})_*\mu_{I,e^{tX}}$ with $B := X/2$, so $\mu_I$ is
$B$-self-decomposable in the sense of \ref{prop:stable-exponent-cone}. $B$ is positive definite, so
$e^{-tB} \to 0$, and $\mu_I$ is rotation invariant and not Gaussian; \ref{prop:anisotropy-budget}(1) gives
$\tr B < (d+2)\lambda_{\min}(B)$, that is $f(X) < 2$. As the largest value of $f$ is attained, $\beta(C) < 2$.
If $\beta(C) \ge 2$, then every family has a Gaussian $\mu_I$. Its image under $x \mapsto \langle e_1, x\rangle$
is a Gaussian law on $\RR$, of some variance $v \ge 0$, so $e^{-g(t^2)} = |\widehat\mu_I(te_1)| = e^{-vt^2/2}$ and
$g(u) = bu$ with $b := v/2$, and $b > 0$ as $g \not\equiv 0$; the end of (2)(b) shows that the
family is the restricted affine Gaussian scale space; these exist, by the converse there, which used only
$C \subseteq \PSD(d)$.

% CHANGED Q-0281 wave 2: entry-free route (the author's decision on W2F-1). The ray order is read
% through prop:stable-order-generated(1) and the definition of the order on PSD(d), in place of
% thm:cone-orders(4).
\emph{(3).} By \ref{prop:stable-order-generated}(1), $\Sigma \preceq_{C_0} \Sigma'$ on $\PD(d)$ means $\Sigma' = c\Sigma$
with $c \ge 1$, so by the definition of $\preceq^+$ the order $\preceq^+_{C_0}$ is the ray order, and
\ref{prop:gw-ray-families}(3) is the description of $\mathfrak F_0$. For $\alpha \le \beta$,
$\preceq^+_{C_\alpha} \subseteq \preceq^+_{C_\beta}$ by (1), and a family restricts; so
$\mathfrak F_\beta \subseteq \mathfrak F_\alpha$. For $0 \le \alpha < \infty$, $C_\alpha$ satisfies the
hypotheses of (2)(c), as its nonzero members are definite by (1), and $\beta(C_\alpha) = \alpha$, since
$f \le \alpha$ on $C_\alpha$ with equality at $\diag(1, \dots, 1, 1+\alpha)$. The clauses of (3) follow,
range by range.
\begin{itemize}
\item $\alpha = 0$: $\mathfrak F_0$ is described above.
\item $0 < \alpha < 2$: the stable clause is (2)(c), and the profile clause is
\ref{prop:similarity-stable-order-families}(2), its order $\le_\alpha$ being $\preceq^+_{C_\alpha}$.
\item $2 \le \alpha < \infty$: (2)(c) with $\beta(C_\alpha) \ge 2$ gives the Gaussians.
\item $\alpha = \infty$: $C_\infty = \PSD(d)$ contains $e_1e_1^{\tp}$, so (2)(b) gives the Gaussians.
\end{itemize}
Strictness for $\alpha < \beta \le 2$: let $\gamma := \alpha$ if $\alpha > 0$ and $\gamma := \beta/2$ if
$\alpha = 0$, so $0 < \gamma < \beta$ and $\gamma \ge \alpha$. By (2)(c) with $\beta(C_\alpha) = \alpha$,
$S_\gamma(I) \in \mathfrak F_\alpha$; and $S_\gamma(I) \notin \mathfrak F_\beta$, by (2)(c) with $\beta(C_\beta) =
\beta > \gamma$ if $\beta < 2$, and because $\mathfrak F_2$ contains only Gaussian laws if $\beta = 2$.
\end{proof}

% CHANGED Q-0281: split from thm:wedge-axis (ADR-0025; records/formalization/PRICING.md): clauses
% (1), (2)(a) and (2)(b), copied verbatim with the parent's clause numbers and the preamble of (2). They
% spend nothing, where the parent's (2)(c) and (3) spend A5 and A3; the parent \uses this node, and its
% text and numbering are unchanged.
% CHANGED Q-0281 wave 2: entry-free route (the author's decision on W2F-1). \uses re-pointed to what
% the proof reads: thm:cone-orders, lem:cone-orders-psd and prop:invariant-orders dropped.
% CHANGED Q-0281 wave 2: blind review W2H-1: the wedge datum that (2) presupposes is the parent's
% paragraph on the orders, which now proves the antisymmetry without ledger A9; \uses gains
% lem:spectral-cones, whose last sentence of (4) that paragraph reads and which spends nothing, and the
% annotation and the proof say so. The Lean is unaffected (IsGWCascadeFamily takes any relation).
% CHANGED Q-0281 wave 2: blind review W2H-2: \uses gains def:cone-orders-curves, the vocabulary of the
% statement; the text is unchanged.
\begin{lemma}[the chain of cones, and the cones that carry only the Gaussian]
\label{lem:wedge-axis-chain}
\uses{def:cone-orders-curves,lem:spectral-cones,prop:stable-order-generated,def:affine-cascade-family,lem:affine-gaussian-transform,def:gw-cascade-family,lem:gw-kernels-from-origin,prop:gw-ray-families,thm:rank-deficient-step,thm:affine-rigidity}
\lean{AffineHemigroup.wedge_axis_chain, AffineHemigroup.wedge_axis_not_psd,
AffineHemigroup.wedge_axis_singular}\leanok
Let $d \ge 2$, and let $C_\alpha$, $\preceq_\alpha$ be as in \ref{prop:stable-order-generated} and
$\preceq^+_C$ as in \ref{thm:cone-orders}(4).
\begin{enumerate}
\item[(1)] \emph{(The chain.)} For $0 \le \alpha < \beta \le \infty$, $C_\alpha \subsetneq C_\beta$ and
$\preceq_\alpha \subsetneq \preceq_\beta$. For $\alpha < \infty$ every nonzero member of $C_\alpha$ is positive
definite, while $C_\infty = \PSD(d)$ contains every nonzero singular positive semidefinite matrix. The union of
the $C_\alpha$, $0 \le \alpha < 2$, is $\{0\} \cup \{X \in \Sym(d) : (d+2)\lambda_{\min}(X) > \tr X\}$, a proper
subset of $C_2$.
\item[(2)] \emph{(Which cones carry general linear families.)} Let $C \subseteq \Sym(d)$ be an $\Orth(d)$-invariant
closed convex cone with $I \in C$ and $C \cap (-C) = \{0\}$, and consider the
$(\GL(d), \preceq^+_C)$-cascade families of kernels (\ref{def:gw-cascade-family}(2)).
\begin{enumerate}
\item[(a)] If $C \not\subseteq \PSD(d)$, there is none.
\item[(b)] If $C \subseteq \PSD(d)$ and $C$ contains a nonzero singular matrix, they are exactly the
restrictions to the pairs $\Sigma \preceq^+_C \Sigma'$ of the affine Gaussian scale spaces
\ref{eq:affine-gaussian}, $b > 0$.
\end{enumerate}
\end{enumerate}
\statusT\quad\emph{(Cites nothing beyond the spectral theorem and the uniqueness of the Fourier
transform: (b) is \ref{thm:rank-deficient-step}. That $(\GL(d), \preceq^+_C)$ is a wedge datum, which (2)
presupposes, is proved in the parent's paragraph on the orders without entry A9, reading only the last
sentence of \ref{lem:spectral-cones}(4), which spends nothing. Clauses (1), (2)(a) and (2)(b) of
\ref{thm:wedge-axis}, split off on 2026-10-04 (Q-0281, ADR-0025); the numbering stays the parent's.)}
\end{lemma}

\begin{proof}
\leanok
% CHANGED Q-0281 wave 2: entry-free route (the author's decision on W2F-1): the parent's proof gains
% the paragraph on the orders, which this part reads.
% CHANGED Q-0281 wave 2: blind review W2H-1: the pointer names the wedge datum of (2).
% CHANGED writing-rules pass (2026-10-14): the parent's paragraph is now headed "The orders".
This is the part of the proof of \ref{thm:wedge-axis} headed the orders, (1), (2)
the kernels from scale zero, (2)(a) and (2)(b). The first of these proves that $\preceq^+_C$ is a partial
order and $(\GL(d), \preceq^+_C)$ a wedge datum, the one whose cascade families (2) considers.
\end{proof}

% not in the draft: round 3 of the external review, report E, section 3.1 and required change 1
% (records/affine/reviews/review-content-review-3-E-section-4.9-5-appendices-D-F-I-claude-fable-5.1.md;
% VERIFY-round3-smaller-groups-paper.md section 1, and VERIFY-round3-landscape-orbits-paper.md section 12
% for report C's laws with a Gaussian part), the author's decision D10 (PLAN-review-response-paper.md
% R3.0, 2026-11-03; option (b) of R3.1 item 1): the converse of the profile clauses of
% prop:similarity-stable-order-families(2) and thm:wedge-axis(3)(c), with any Gaussian part, the same
% along a cone order, and with the polar criterion (ledger A2) the whole cell. A new node, so that the
% two green nodes keep a Lean that covers all of them (ADR-0025); they point here in their annotations.
% What writing it down found: (i) the "open converse of prop:orbit-step-cone" that those annotations
% named is the single-step converse; for a cascade family it is not needed, the increments being
% infinitely divisible (thm:cascade-infinitely-divisible(2)) and an infinitely divisible increment
% ordering the Levy densities (prop:similarity-orbit-family(1)); (ii) the cone version needs no path
% condition for the whole order, only the chain s -> e^{stX} for one X attaining beta(C); (iii) the
% whole cell applies the polar criterion (A2) and A3(i) directly to mu_I, already infinitely
% divisible, so entry A4, which the forward direction of prop:isotropic-polar-reading(1) spends to
% make a self-decomposable law infinitely divisible, is not needed (round 3: the proof first read
% that node, whose uses-edge brought A4 into the trust base, and was rewritten to the direct route,
% with the sphere mean in place of its radial-image uniqueness; the edge and def:isotropic-origin-datum
% left the uses list); (iv) right-continuity is the hypothesis
% of prop:similarity-orbit-family(1), and the whole cell takes the right-continuous representative of
% the polar profile. Priced at writing, beside the statement: clause (1), an assembly of
% stable_order_gl_row, gw_kernels_from_origin_converse (the Gaussian factor), chain_infinitelyDivisible,
% sim_orbit_id_increment at h = id and C(Sigma) = a Sigma, wedge_axis_index and dilRatio_exponents,
% about 120-180 Lean lines; clause (2), a form of polar_reading_forward that takes infinite
% divisibility as a hypothesis (about 20 lines), the right-continuous modification of a monotone
% profile and the identification of the laws, about 150-250 lines. The \lean tag names the two
% declarations to be written in the desktop session; until then the node is \notready, proved in print.
% CHANGED the author's decision D14 (2026-11-04; PLAN-review-response-paper.md R3.0, the round-3
% referee's finding 1): a widening, the author's, not a safe-direction change. The profile k of the
% statement's preamble, which clause (1) reads, is no longer assumed right-continuous; clause (2) keeps
% "right-continuous" on its own bound k, where it picks one profile for each law (without it the set of
% laws would be the same). The proof gains the paragraph "(1), the right-continuous modification": the
% modification k(r+) has the same law and the same lower exponent, so the necessity may take k
% right-continuous, as prop:similarity-orbit-family(1) asks. Remark (iv) above now concerns the
% necessity of (1) and the cell. The "Not claimed" sentence on profiles that are not right-continuous is
% dropped. Price added to clause (1): the modification is a profile with the same law and the same
% lower exponent, about 60-100 Lean lines, shared with clause (2)'s modification of the polar profile.
% CHANGED round 3, the node's Lean: the two clauses are proved, in WedgeAxisProfiles.lean (490 lines,
% priced 240-380 with the modification) and WedgeAxisCell.lean (479 lines, priced 170-270), headers
% included. The tag names the four declarations that state the node and rowLaws_coe_eq_coneLaws, for
% F_alpha = F(C_alpha). What writing it down found: (i) the chain lemma asks for a partial order, and
% for C inside PSD(d) the antisymmetry of the cone order is read from the Loewner order, so A9 stays
% unspent; (ii) the Gaussian part goes through one orbit, of N(0, aI) * nu_I, so that (a)-(c) of the
% kernels from the origin are those of an orbit of a rotation-invariant law; (iii) the guard prints
% A3(iii) in the necessity half too, so the Assignment says so; (iv) the hypotheses alpha < 2 of the
% cell and beta(C) < 2 of the cones are not read: the form of mu_I and clause (1) exclude every profile
% once the number is at least 2, since beta_-(k) < 2. So those Gaussian ends hold on A2, A3 and A6, where
% the Not claimed sentence reaches them through A5. The statement keeps the hypotheses: dropping them
% would be a widening, the author's. Route: the proof's "(2), the form of mu_I" no longer proves the mean
% k right-continuous by monotone convergence; the Lean takes the right-continuous modification of the
% first paragraph on it, so the proof says that. Nothing else in the proof moved: the Lean's packaging
% of (ii) and (i) reads the same conditions, which the printed sufficiency and first paragraph also prove.
% CHANGED the author's decision D16 (2026-11-04; PLAN-review-response-paper.md R3.0): a widening, the
% author's. Clause (2) drops alpha < 2 and beta(C) < 2, which neither proof read (item (iv) above). The
% cell now runs over 0 < alpha < infinity, and the statement says what is left from alpha = 2 on, as
% beta_-(k) < 2; so the sets of laws of thm:wedge-axis(2)(c) and (3)(d), for finite alpha, also hold on A6,
% A3 and A2, without A5, and the end of its (2)(b), which spends nothing, then gives the families. That
% theorem's own proof is not rerouted: prop:anisotropy-budget(1) spends A5 in any case, so the entry stays
% in the paper, and A2 and A6 are no lighter a route. Two declarations for the Gaussian ends. The blind
% referee (records/affine/REFEREE-lean-round3.md, "D16") corrected the first wording of these sentences.
\begin{proposition}[the general linear row along the chain: the profiles, and the whole cell]
\label{prop:wedge-axis-profiles}
\uses{thm:wedge-axis,def:gw-cascade-family,lem:gw-kernels-from-origin,thm:cascade-infinitely-divisible,prop:stable-order-generated,prop:similarity-stable-order-families,lem:stable-order-psd,prop:orbit-step-cone,lem:dilation-ratio,prop:similarity-orbit-family,prop:gw-ray-families,thm:affine-rigidity,thm:cited-A2,thm:cited-A3,thm:cited-A6}
\lean{AffineHemigroup.wedge_axis_profiles, AffineHemigroup.wedge_axis_profiles_stable,
AffineHemigroup.rowLaws_coe_eq_coneLaws, AffineHemigroup.wedge_axis_cell,
AffineHemigroup.wedge_axis_cell_gaussian, AffineHemigroup.wedge_axis_cone_laws,
AffineHemigroup.wedge_axis_cone_laws_gaussian}\leanok
Let $d \ge 2$, and let $C_\alpha$, $\preceq^+_C$, the number $\beta(C)$ and $\mathfrak F_\alpha$ be as in
\ref{thm:wedge-axis}. For a cone $C$ as in \ref{thm:wedge-axis}(2) with $C \subseteq \PSD(d)$ and no nonzero
singular member, the case (2)(c) there, write $\mathfrak F(C)$ for the set of laws $\mu_I$ of
$(\GL(d), \preceq^+_C)$-cascade families of kernels, so that $\mathfrak F_\alpha = \mathfrak F(C_\alpha)$ for
$0 \le \alpha < \infty$. Let $k$ be a profile as in \ref{prop:orbit-step-cone}, with dilation
ratio $R$ and lower exponent $\beta_-(k)$ (\ref{lem:dilation-ratio}(3)), let $\nu_I$ be its law
(\ref{prop:orbit-step-cone}(1)), and for $a \ge 0$ write $N(0, aI)$ for the centred Gaussian law with
covariance $aI$.
\begin{enumerate}
\item \emph{(The profiles.)} Let $C$ be such a cone, with $\beta(C) > 0$. Then $N(0, aI) * \nu_I \in \mathfrak F(C)$
if and only if $\beta(C) \le \beta_-(k)$. In particular, for $0 < \alpha < 2$, $N(0, aI) * \nu_I \in \mathfrak F_\alpha$
if and only if $R(\lambda) \ge \lambda^\alpha$ for every $\lambda \ge 1$, and then
$\Sigma \mapsto N(0, a\Sigma) * (\Sigma^{1/2})_*\nu_I$, $\Sigma \in \PSD(d)$, is the family of kernels from scale
zero of a $(\GL(d), \preceq^+_{C_\alpha})$-cascade family of kernels.
\item \emph{(The whole cell.)} For $0 < \alpha < \infty$, $\mathfrak F_\alpha$ consists of the laws $N(0, aI)$, $a > 0$,
and the laws $N(0, aI) * \nu_I$, $a \ge 0$, of the right-continuous profiles $k$ with $\beta_-(k) \ge \alpha$.
From $\alpha = 2$ on there is no such profile, so $\mathfrak F_\alpha$ consists of the laws $N(0, aI)$, $a > 0$.
For a cone $C$ as above, $\mathfrak F(C) = \mathfrak F_{\beta(C)}$: the laws $\mu_I$ of the families
along $\preceq^+_C$ depend on $C$ only through $\beta(C)$.
\end{enumerate}
\statusA\quad\ledger{A6}\quad\ledger{A3}\quad\ledger{A2}\quad\emph{(\textbf{Assignment.} \ledger{A6},
Khintchine's theorem, through \ref{thm:cascade-infinitely-divisible}(2): in (1) for the infinitely divisible
increment of one step joined by a chain, and in (2) for the infinite divisibility of $\mu_I$.
\ledger{A3}(i)--(ii) through \ref{prop:similarity-orbit-family}(1), in the necessity half of (1), which reads
the ordering of the \Levy{} densities off an infinitely divisible increment; \ledger{A3}(iii) through
\ref{prop:orbit-step-cone}(1) and \ref{prop:similarity-stable-order-families}(2), in the sufficiency half, for
the law of the profile and its increments, and through \ref{prop:similarity-orbit-family}(1), in the necessity
half, which reads the rotation invariance of that law from \ref{prop:orbit-step-cone}(1); \ledger{A3}(i) again in
(2), for the generating triplet of $\mu_I$, and \ledger{A3}(iii) for the law of the profile it produces.
\ledger{A2}, the polar criterion, in (2) only, applied directly to $\mu_I$, which is infinitely divisible by
\ledger{A6}. The isotropic paper's polar reading (its Proposition~3.11) is not read: its forward direction
spends entry A4 to make a self-decomposable law infinitely divisible, and here $\mu_I$ is so already, so A4
is not needed. Of \ref{thm:wedge-axis} the node reads only what spends nothing:
the number $\beta(C)$ and $C \subseteq C_{\beta(C)}$ in (2)(c), the paragraph on the orders of its proof
(the orders, their causal curves, the restriction of a family to a smaller order), and
$\beta(C_\alpha) = \alpha$; not the anisotropy budget, entry A5. \textbf{What it is worth.} It completes the
profile clause of \ref{thm:wedge-axis}(3)(c) and the converse of \ref{prop:similarity-stable-order-families}(2):
along the stable order of index $\alpha$ the orbits of the profiles that are general linear families are
exactly those whose lower exponent is at least $\alpha$, with any Gaussian part, and with the polar
criterion the cell $(\GL(d), \le_\alpha)$ is described completely. Along a cone order the same holds with
$\beta(C)$ in place of $\alpha$, so two cones with the same $\beta(C)$ carry the same laws $\mu_I$. From
$\alpha = 2$ on no profile qualifies, as $\beta_-(k) < 2$, so (2) gives the sets of laws of
\ref{thm:wedge-axis}(2)(c) at $\beta(C) \ge 2$ and of (3)(d) for finite $\alpha$ on \ledger{A6}, \ledger{A3} and
\ledger{A2}, without entry A5, and the end of (2)(b) of that theorem, which spends nothing, then gives the
families. That theorem keeps its own route through the anisotropy budget; the Lean states the laws here. The
single-step converse of \ref{prop:orbit-step-cone} is not used: a step of a cascade has an infinitely
divisible increment, and only such an increment is known to order the densities (a positive increment
need not be infinitely divisible: the Gaussian--\Matern{} step in the annotation of
\ref{prop:similarity-orbit-family}). For $C = C_\alpha$ the step $e^{tX}$ of the proof, with
$X = \diag(1, \dots, 1, 1 + \alpha)$, is the step of the proof of \ref{lem:dilation-ratio}(4).
\textbf{Not claimed}: the single-step converse; any cell of a smaller group; the cones of
\ref{thm:wedge-axis}(2)(a)--(b).
\textbf{Right-continuity.} Clause (1) holds for every profile: the right-continuous modification
$r \mapsto k(r{+})$ has the same law and the same lower exponent (the proof), and the necessity reads
\ref{prop:similarity-orbit-family}(1), which asks for a right-continuous profile, only for the modification.
Clause (2) keeps the word: without it the set of laws is the same, by the same step, and with it each law
$N(0, aI) * \nu_I$ has one profile, its \Levy{} measure fixing $k$ almost everywhere and right-continuity
then everywhere.)}
\end{proposition}

\begin{proof}
\leanok
\emph{The orders.} For $0 < \gamma < 2$, $\preceq^+_{C_\gamma}$ is the order $\le_\gamma$ of
\ref{prop:similarity-stable-order-families}(1): on $\PD(d)$ both are $\preceq_\gamma$
(\ref{prop:stable-order-generated}(1)), and both add the origin and the rays. Let $C$ be a cone as in the
statement. By \ref{thm:wedge-axis}(2)(c) the largest value $\beta(C)$ of $\tr X/\lambda_{\min}(X) - d$ over
$X \in C \setminus \{0\}$ is attained and $C \subseteq C_{\beta(C)}$. So $\preceq^+_C$ is contained in
$\preceq^+_{C_{\beta(C)}}$, a $C$-causal curve being $C_{\beta(C)}$-causal, and a
$(\GL(d), \preceq^+_{C_{\beta(C)}})$-cascade family restricts to a $(\GL(d), \preceq^+_C)$-cascade family with the
same $\mu_I$ (the proof of \ref{thm:wedge-axis}, the orders). For $X \in C$ and $t \ge 0$ the curve
$\zeta(s) := e^{stX}$, $s \in [0,1]$, is a chain from $I$ to $e^{tX}$ for $\preceq^+_C$ as in
\ref{thm:cascade-infinitely-divisible}(2). It is continuous, and for $s \le s'$, $I \preceq^+_C e^{(s'-s)tX}$
as $(s'-s)tX \in C$ (the proof of \ref{thm:wedge-axis}, the orders). The matrix $A := e^{stX/2}$ carries
$I$ to $\zeta(s)$ and $e^{(s'-s)tX}$ to $\zeta(s')$, so $\zeta(s) \preceq^+_C \zeta(s')$ by the
$\GL(d)$-invariance of the order.

\emph{(1), sufficiency.} Let $\beta(C) \le \beta_-(k)$ and put $\gamma := \beta(C)$. Then $0 < \gamma < 2$, as
$\beta_-(k) < 2$, and $R(\lambda) \ge \lambda^\gamma$ for every $\lambda \ge 1$ (\ref{lem:dilation-ratio}(3)). By
\ref{prop:similarity-stable-order-families}(2) the laws $\nu_\Sigma := (\Sigma^{1/2})_*\nu_I$, $\Sigma \in \PSD(d)$,
are the kernels from scale zero of a $(\GL(d), \le_\gamma)$-cascade family $(\nu_{\Sigma,\Sigma'})$. Put
$\mu_\Sigma := N(0, a\Sigma) * \nu_\Sigma$. Its transform $e^{-a\langle\omega,\Sigma\omega\rangle/2}\,\widehat\nu_\Sigma(\omega)$
has no zero, as $\widehat\nu_\Sigma(\omega) = e^{-\psi_I(\Sigma^{1/2}\omega)}$ (\ref{prop:orbit-step-cone}(1)). The
conditions (a)--(d) of \ref{lem:gw-kernels-from-origin}(2) hold for $G = \GL(d)$ and the order $\le_\gamma$, as
they hold for each factor.
\begin{itemize}
\item (a): $\mu_0 = \delta_0$, and $\mu_I \ne \delta_0$, as $\widehat\mu_I = e^{-a|\omega|^2/2 - \psi_I}$ and $\psi_I$
is not identically $0$, $k$ being positive near $0$.
\item (b): $A_*N(0, a\Sigma) = N(0, aA\Sigma A^{\tp})$ and $A_*\nu_\Sigma = \nu_{A\cdot\Sigma}$ for $A \in \GL(d)$.
\item (c): both transforms are continuous in $\Sigma$.
\item (d): for $\Sigma \le_\gamma \Sigma'$, $\Sigma' - \Sigma \ge 0$ (\ref{prop:similarity-stable-order-families}(1)), so
$\rho := N(0, a(\Sigma' - \Sigma)) * \nu_{\Sigma,\Sigma'}$ is a probability measure, and $\mu_{\Sigma'} = \mu_\Sigma * \rho$.
\end{itemize}
By \ref{lem:gw-kernels-from-origin}(3) the $\mu_\Sigma$ are the kernels from scale zero of a
$(\GL(d), \le_\gamma)$-cascade family, whose law at $I$ is $N(0, aI) * \nu_I$, and it restricts to $\preceq^+_C$.

\emph{(1), the right-continuous modification.} Put $\tilde k(r) := k(r{+}) = \lim_{s \downarrow r}k(s)$ for $r > 0$.
It is nonincreasing and right-continuous, and $\tilde k \le k$, so $\tilde k$ satisfies the
integrability conditions that $k$ does. If $k(r_0) > 0$, then $\tilde k \ge k(r_0)$ on
$(0, r_0)$, so $\tilde k$ is a profile as in \ref{prop:orbit-step-cone}. A monotone function has at most
countably many discontinuities, so $\tilde k = k$ off a countable subset of $(0, \infty)$, and the densities
$k(|x|)|x|^{-d}$ and $\tilde k(|x|)|x|^{-d}$ agree off the set of the $x$ with $|x|$ in that countable set,
a null set of $\RR^d$. So the two profiles have the same
$\psi_I$ and the same law $\nu_I$. They also have the same lower exponent. Let $\tilde R$ be the dilation ratio of
$\tilde k$ and $\lambda \ge 1$. Every $s > 0$ has $k(s/\lambda) \ge R(\lambda)k(s)$, by the definition of $R$ when
$k(s) > 0$ and trivially otherwise. Letting $s$ decrease to $r$ gives $\tilde k(r/\lambda) \ge R(\lambda)\tilde k(r)$,
so $\tilde R(\lambda) \ge R(\lambda)$. Conversely, let $\alpha \ge 0$ with $\tilde R(\tau) \ge \tau^\alpha$ for every
$\tau \ge 1$, let $\lambda > 1$, and let $k(r) > 0$. For $r/\lambda < s < r$, $\tilde k(s) \ge k(r) > 0$, as
$k \ge k(r)$ on $(s, r)$, and $\tau := s\lambda/r > 1$. So
$k(r/\lambda) \ge \tilde k(r/\lambda) = \tilde k(s/\tau) \ge \tau^\alpha\tilde k(s) \ge (s\lambda/r)^\alpha k(r)$, and
letting $s$ increase to $r$ gives $k(r/\lambda) \ge \lambda^\alpha k(r)$. So $R(\lambda) \ge \lambda^\alpha$, which
holds at $\lambda = 1$ as well. By \ref{lem:dilation-ratio}(3), applied to $k$ and to $\tilde k$, the numbers
$\alpha \ge 0$ with $\alpha \le \beta_-(k)$ are those with $\alpha \le \beta_-(\tilde k)$, and
$\beta_-(k) = \beta_-(\tilde k)$. Both sides of the equivalence in (1) are therefore the same for $k$ and $\tilde k$.
In the necessity $k$ may be replaced by $\tilde k$, and so taken right-continuous.

\emph{(1), necessity.} By the modification we may take $k$ right-continuous. Let $\mu_I := N(0, aI) * \nu_I$ be the law at $I$ of a $(\GL(d), \preceq^+_C)$-cascade
family $(\mu_{\Sigma,\Sigma'})$. For $\Sigma \in \PD(d)$, (K3) with $A = \Sigma^{1/2}$ at the pair $(0, I)$ gives
$\mu_\Sigma = (\Sigma^{1/2})_*\mu_I = N(0, a\Sigma) * (\Sigma^{1/2})_*\nu_I$. Let $X \in C \setminus \{0\}$ attain
$\beta(C)$, so that $\tr X = (d + \beta(C))\lambda$ with $\lambda := \lambda_{\min}(X) > 0$, $X$ being positive definite. Let $t > 0$ and $M := e^{tX}$. By the chain of the first paragraph and
\ref{thm:cascade-infinitely-divisible}(2) the increment $\rho := \mu_{I,M}$ is infinitely divisible, and (K1) at
$0 \preceq^+_C I \preceq^+_C M$ gives $\mu_M = \mu_I * \rho$. So \ref{prop:similarity-orbit-family}(1), with $h$ the
identity and the Gaussian covariance $\Sigma \mapsto a\Sigma$, puts the step $I \to M$ in the step set of
\ref{prop:orbit-step-cone}(2): $\det M \le \lambda_{\min}(M)^dR(\lambda_{\min}(M)^{1/2})^2$. Here
$\lambda_{\min}(M) = e^{t\lambda}$ and $\det M = e^{t\tr X} = e^{t(d+\beta(C))\lambda}$, so with $u := e^{t\lambda/2}$ the
inequality reads $u^{2(d+\beta(C))} \le u^{2d}R(u)^2$, that is $R(u) \ge u^{\beta(C)}$. As $t$ ranges over
$(0, \infty)$, $u$ ranges over $(1, \infty)$, and $R(1) = 1$. So $\beta(C) \le \beta_-(k)$ by
\ref{lem:dilation-ratio}(3). For $C = C_\alpha$ one may take $X = \diag(1, \dots, 1, 1 + \alpha)$, and $M$ is the
step of the proof of \ref{lem:dilation-ratio}(4).

\emph{(1), the stable orders.} For $0 < \alpha < 2$, $C_\alpha$ is a cone as in the statement with
$\beta(C_\alpha) = \alpha$ (the proof of \ref{thm:wedge-axis}, (3)), and $\mathfrak F_\alpha = \mathfrak F(C_\alpha)$. By
\ref{lem:dilation-ratio}(3), $\alpha \le \beta_-(k)$ if and only if $R(\lambda) \ge \lambda^\alpha$ for every
$\lambda \ge 1$, and the family is the one built in the sufficiency with $\gamma = \alpha$, along
$\le_\alpha = \preceq^+_{C_\alpha}$.

\emph{(2), the form of $\mu_I$.} Let $C$ be a cone as in the statement and $\mu_I \in \mathfrak F(C)$. The ray order
is contained in $\preceq^+_C$, so the family restricts to a $(\GL(d), \text{rays})$-cascade family with the same
$\mu_I$, and by \ref{prop:gw-ray-families}(3) and (1), $\mu_I$ is rotation invariant, self-decomposable and not
$\delta_0$, and $\widehat\mu_I$ has no zero. The pair $0 \preceq^+_C I$ is joined by the chain $s \mapsto sI$, so
$\mu_I = \mu_{0,I}$ is infinitely divisible (\ref{thm:cascade-infinitely-divisible}(2)). Its transform is radial,
and real as $\mu_I$ is invariant under $-I$. Having no zero, it is $\widehat\mu_I(\omega) = e^{-F(|\omega|)}$ with
$F(r) := -\log\widehat\mu_I(re_1) \ge 0$. By \ledger{A3}(i), $\mu_I$ has a generating triplet $(A, \nu, \gamma)$.
The real part of $e^{i\theta} - 1 - i\theta\mathbf 1_{\{|x| \le 1\}}$ is $\cos\theta - 1$, so taking moduli in the
representation,
\[
  F(|\omega|) = \tfrac12\langle\omega, A\omega\rangle + \int_{\RR^d}\bigl(1 - \cos\langle\omega,x\rangle\bigr)\,\nu(dx) ,
  \qquad \omega \in \RR^d .
\]
As $\mu_I$ is self-decomposable, the polar criterion \ledger{A2} gives
$\nu(B) = \int_S\lambda(d\xi)\int_0^\infty\mathbf 1_B(r\xi)\,k_\xi(r)\,dr/r$ for Borel $B$, with $\lambda$ a finite
measure on $S = S^{d-1}$ and $k_\xi(r) \ge 0$ measurable in $\xi$ and nonincreasing in $r$. A nonincreasing $k_\xi$
differs from its right limit $k_\xi(r+) = \sup\{k_\xi(q) : q \in \mathbb Q,\ q > r\}$ at countably many $r$ only,
so replacing $k_\xi(r)$ by $k_\xi(r+)$ changes no inner integral, and it makes the kernel jointly measurable in
$(\xi, r)$, as a countable supremum of measurable functions. Let $\sigma_{d-1}$ be the surface measure of $S$, with total
$|S^{d-1}| = \sigma_{d-1}(S)$, and put
\[
  k(r) := \frac{1}{|S^{d-1}|}\int_S k_\xi(r+)\,\lambda(d\xi) \in [0, \infty] , \qquad r > 0 .
\]
It is nonincreasing. Exchanging the integrals (the integrands are nonnegative) and passing to polar coordinates
$x = r\eta$, $dx = r^{d-1}\,dr\,\sigma_{d-1}(d\eta)$,
\[
  \int_{\RR^d}(|x|^2 \wedge 1)\,\nu(dx) = |S^{d-1}|\int_0^\infty(r^2 \wedge 1)\,k(r)\,\frac{dr}{r}
  = \int_{\RR^d}(|x|^2 \wedge 1)\,k(|x|)\,|x|^{-d}\,dx ,
\]
which is finite by \ledger{A3}(i). So $k$ is finite: a nonincreasing function infinite at some $r_0$ is infinite
on $(0, r_0]$, and the integral would diverge. The left side of the display for $F$ depends on $\omega$ only
through $|\omega|$, so it equals the mean of the right side at the points $|\omega|\theta$, $\theta \in S$, against
$\sigma_{d-1}/|S^{d-1}|$. The Gaussian term has the mean $\tfrac12 a|\omega|^2$ with $a := \tr A/d \ge 0$, as
$\int_S\theta_i\theta_j\,\sigma_{d-1}(d\theta) = \delta_{ij}|S^{d-1}|/d$. For the other term let
$c(s) := |S^{d-1}|^{-1}\int_S(1 - \cos(s\langle\theta,\xi\rangle))\,\sigma_{d-1}(d\theta)$ for $s \ge 0$ and a unit vector $\xi$. It does not depend on $\xi$, since $\sigma_{d-1}$ is invariant under a rotation carrying one unit vector to
another. Exchanging the integrals as before and then passing to polar coordinates, with
$c(r|\omega|) = |S^{d-1}|^{-1}\int_S(1 - \cos\langle\omega, r\eta\rangle)\,\sigma_{d-1}(d\eta)$ (take $\xi = \omega/|\omega|$
when $\omega \ne 0$),
\[
  F(|\omega|) = \tfrac12 a|\omega|^2 + |S^{d-1}|\int_0^\infty c(r|\omega|)\,k(r)\,\frac{dr}{r}
  = \tfrac12 a|\omega|^2 + \int_{\RR^d}\bigl(1 - \cos\langle\omega,x\rangle\bigr)\,k(|x|)\,|x|^{-d}\,dx ,
  \qquad \omega \in \RR^d .
\]
If $k \equiv 0$, then $\widehat\mu_I(\omega) = e^{-a|\omega|^2/2}$ and
$\mu_I = N(0, aI)$, with $a > 0$ as $\mu_I \ne \delta_0$. Otherwise $k$ is a profile as in
\ref{prop:orbit-step-cone}, and so is its right-continuous modification $\tilde k$, with the same exponent (the paragraph
on the modification). So $\widehat\mu_I = e^{-a|\omega|^2/2}\,\widehat\nu_I$ for the law $\nu_I$ of $\tilde k$, and
$\mu_I = N(0, aI) * \nu_I$ by the uniqueness of the transform.

\emph{(2), the equality.} If $\beta(C) = 0$, then $C \subseteq C_0 = \{cI : c \ge 0\}$
(\ref{prop:stable-order-generated}(1)), and $C_0 \subseteq C$ as $I \in C$. So $C = C_0$, and there is nothing to prove. Let $\beta(C) > 0$. By the first paragraph, $\mathfrak F_{\beta(C)} \subseteq \mathfrak F(C)$.
Conversely let $\mu_I \in \mathfrak F(C)$, of one of the two forms above. A law $N(0, aI)$, $a > 0$, is the law
at $I$ of the affine Gaussian scale space with constant $a/2$, an affine cascade family
(\ref{thm:affine-rigidity}), and so of its restriction to the pairs of $\preceq^+_{C_{\beta(C)}}$, which are
Loewner pairs, $C_{\beta(C)}$ being contained in $\PSD(d)$ (the proof of \ref{thm:wedge-axis}, the orders). So it lies in $\mathfrak F_{\beta(C)}$. A law $N(0, aI) * \nu_I$ has $\beta(C) \le \beta_-(k)$ by (1) for $C$, and so
lies in $\mathfrak F(C_{\beta(C)}) = \mathfrak F_{\beta(C)}$ by (1) for $C_{\beta(C)}$, whose number is $\beta(C)$. Hence
$\mathfrak F(C) = \mathfrak F_{\beta(C)}$. For $C = C_\alpha$, $0 < \alpha < \infty$, a cone as in the statement with
$\beta(C_\alpha) = \alpha$ (the proof of \ref{thm:wedge-axis}, (3)), the same two cases describe
$\mathfrak F_\alpha$: each member has one of the two forms, each $N(0, aI)$ with $a > 0$ lies in it, and
$N(0, aI) * \nu_I$ lies in it exactly when $\beta_-(k) \ge \alpha$, by (1). For $\alpha \ge 2$ the second case does
not occur, as $\beta_-(k) < 2$ for every profile (\ref{lem:dilation-ratio}(3)).
\end{proof}
